\chapter{Electromyography and Nerve-Conduction Study (EMG/NCS)}

\section*{What Is This Test?}
Electromyography and nerve-conduction studies evaluate the electrical function of muscles and peripheral nerves. Nerve-conduction testing measures how quickly and strongly nerves transmit signals, while needle EMG records electrical activity in selected muscles at rest and during contraction. They are commonly performed together.

\section*{Alternative Names}
EMG; Electromyogram; NCS; Nerve-Conduction Study; Nerve-Conduction Velocity; Electrodiagnostic Study

\section*{Why It Is Ordered}
Evaluation of weakness, numbness, tingling, muscle wasting, cramps, or radiating limb pain; suspected carpal tunnel syndrome, radiculopathy, peripheral neuropathy, nerve injury, motor-neuron disease, or a primary muscle disorder; and localization of nerve or muscle dysfunction.

\section*{How This Test Is Performed}
Surface electrodes deliver brief, low-level electrical impulses to a nerve while responses are recorded. For the needle portion, a very fine electrode is inserted into selected muscles and activity is recorded during relaxation and contraction. The study usually takes 30--90 minutes, depending on the number of limbs and muscles examined.

\section*{Preparation, Risks, and Limitations}
Wear loose clothing and keep the skin free of lotions. Tell the clinician about a pacemaker or defibrillator, anticoagulants, bleeding disorders, or infection risk. Nerve stimulation may feel like brief tingling; needle testing can cause temporary soreness or bruising. Results depend on timing, disease stage, temperature, and the muscles and nerves selected, so they must be interpreted with the examination and history.

\section*{References}
\begin{enumerate}
  \item MedlinePlus. Electromyography (EMG) and Nerve-Conduction Studies.
  \item American Association of Neuromuscular \& Electrodiagnostic Medicine. Electrodiagnostic testing.
\end{enumerate}

\clearpage
\section*{Regulatory Status and Authorization Date}
Regulatory status applies to the specific assay, instrument, manufacturer, and intended use rather than to the test name alone. No single approval date can be assigned to this test category. For a commercial assay, verify the current product in the relevant regulator's database: the U.S. Food and Drug Administration (FDA) clearance, approval, or authorization; India's Central Drugs Standard Control Organisation (CDSCO) licence or permission; the European Union In Vitro Diagnostic Regulation (EU IVDR) CE certificate issued through the applicable conformity-assessment route; or the United Kingdom Medicines and Healthcare products Regulatory Agency (MHRA) registration and UKCA/CE status. Laboratory-developed tests may not have an individual FDA, CDSCO, EU IVDR, or MHRA approval date.
\clearpage
